\chapter{استنتاج‌ها و اثبات‌های ریاضی}
\label{app:math_derivations}

این پیوست استنتاج‌های ریاضی پشتیبان روش‌شناسی را ارائه می‌کند.

%==============================================================================
\section{گرادیان تابع زیان نمونه‌اولیه‌ای}
%==============================================================================

برای شبکه‌های نمونه‌اولیه‌ای با تابع زیان آنتروپی متقاطع، گرادیان زیان نسبت به بردار نهفتۀ پرس‌وجو $f_\theta(x)$ به‌صورت زیر استخراج می‌شود. فرض کنید $p_k = \exp(z_k) / \sum_{k'} \exp(z_{k'})$ که در آن $z_k = \langle f_\theta(x), \mathbf{c}_k \rangle / \tau$. برای کلاس درست $y$:
\begin{equation}
\frac{\partial \mathcal{L}}{\partial f_\theta(x)} = \frac{1}{\tau} \sum_k (p_k - \mathbb{1}[y=k]) \mathbf{c}_k
\end{equation}
گرادیان بردار نهفتۀ پرس‌وجو را به سمت نمونه‌اولیه‌ی درست می‌راند و از نمونه‌اولیه‌ی نادرست دور می‌کند، با بزرگی مقیاس‌شده با دما $\tau$.

%==============================================================================
\section{\texorpdfstring{فرمول‌بندی زیان کانونی }{فرمول‌بندی زیان کانونی (Focal Loss)}}
%==============================================================================

\lr{Focal loss}~\cite{lin2017focal} عبارت است از:
\begin{equation}
\mathcal{L}_{\text{\lr{focal}}} = -\alpha (1-p_t)^\gamma \log(p_t)
\end{equation}
که در آن $p_t$ احتمال پیش‌بینی‌شده برای کلاس درست است. برای $\gamma = 0$، این به آنتروپی متقاطع کاهش می‌یابد. برای $\gamma > 0$، نمونه‌های آسان ($p_t$ بالا) با $(1-p_t)^\gamma$ کم‌اهمیت می‌شوند و زیان روی نمونه‌های سخت تمرکز می‌کند. ما $\gamma = 2$ و $\alpha$ از وزن‌های کلاس برای نامتعادل \pnum{42}/\pnum{58} استفاده می‌کنیم.

%==============================================================================
\section{آزمون \lr{DeLong} برای مقایسه‌ی \lr{ROC}}
%==============================================================================

آزمون \lr{DeLong} دو منحنی \lr{ROC} همبسته را با محاسبه‌ی کوواریانس برآوردهای \lr{AUC} مقایسه می‌کند. تحت فرض صفر که دو منحنی \lr{AUC} برابر دارند، آماره‌ی آزمون به‌طور مجانبی نرمال است. این آزمون فقط زمانی قابل اتکاست که دو مدل روی همان نمونه‌های آزمون و با بردارهای احتمال هم‌تراز ارزیابی شده باشند. بنابراین در نسخۀ نهایی پایان‌نامه، مقایسه‌های بین‌پروتکلی توصیفی گزارش می‌شوند و آزمون جفتی کامل \lr{DeLong}، \lr{McNemar} یا بوت‌استرپ جفتی به‌عنوان کار آینده توصیه می‌شود.

%==============================================================================
\section{بازه‌های اطمینان \lr{bootstrap}}
%==============================================================================

برای بازه‌های اطمینان \lr{bootstrap}، مجموعه‌ی آزمون را \pnum{1,000} بار با جایگزینی زیرنمونه می‌کنیم، برای هر زیرنمونه \lr{AUROC} محاسبه می‌کنیم، و صدک‌های \pnum{2.5} و \pnum{97.5} را به‌عنوان بازه‌ی اطمینان \pnum{95}\% می‌گیریم. این فرض می‌کند مجموعه‌ی آزمون نماینده است؛ برای مجموعه‌های آزمون کوچک، بازه ممکن است پهن باشد. مجموعه‌ی آزمون \pnum{400} نمونه‌ای ما پهنای بازه‌ی حدود \prange{0.10}{0.12} برای \lr{AUROC} می‌دهد.

%==============================================================================
\section{استنتاج پراکسی \lr{SOFA}}
%==============================================================================

نمره‌ی \lr{SOFA} (\lr{Sequential Organ Failure Assessment}) ابزار بالینی برای سنجش نارسایی اندام است. از پراکسی ساده‌شده استفاده می‌کنیم:
\begin{equation}
\text{\lr{SOFA}}_{\text{\lr{proxy}}} = w_1(1 - \text{\lr{creatinine}}) + w_2(1 - \text{\lr{platelets}}) + w_3(1 - \text{\lr{SpO}}_2) + w_4(1 - \text{\lr{lactate}})
\end{equation}
که در آن هر جزء به $[0,1]$ نرمال شده و $w_i$ وزن‌های برابر (\pnum{0.25}) هستند. مقادیر بالاتر عملکرد بدتر را نشان می‌دهند. از \lr{(1 - value)} استفاده می‌کنیم تا کراتینین بالا (بد) به \lr{(1 - creatinine)} پایین نگاشت شود. پراکسی با امکان‌پذیری همبسته است: \lr{SOFA} بالاتر $\to$ امکان‌پذیری پایین‌تر، مطابق انتظار بالینی.

%==============================================================================
\section{شاخص شوک}
%==============================================================================

شاخص شوک = ضربان قلب / فشار خون سیستولیک. بازه‌ی نرمال: \pnum{0.5} تا \pnum{0.7}. مقادیری برای شاخص شوک فراتر از \pnum{1.0} ناپایداری همودینامیک را نشان می‌دهند. با \lr{clip} و \lr{min-max} به $[0,1]$ نرمال می‌کنیم. شاخص شوک بالا با ترخیص ناممکن (ناپایداری همودینامیک تأخیر در ترخیص ایمن) همراه است.
